% =====================================================================
%  クラウドコンピューティング(6) 課題レポート
%  コンパイル: lualatex report.tex （app5_ver3.js と同じ階層に置くこと）
% =====================================================================
\documentclass[11pt,a4paper]{ltjsarticle}

\usepackage{luatexja-fontspec}
\usepackage[margin=25mm]{geometry}
\usepackage[hidelinks]{hyperref}
\usepackage{listings}
\usepackage{xcolor}

% ---- ここだけ自分の情報に書き換えてください -------------------------
\newcommand{\studentID}{（25G1059）}
\newcommand{\studentName}{佐藤 涼菜}
\newcommand{\siteURL}{http://4.217.216.124:8080/}
% ---------------------------------------------------------------------

% ---- ソースコード掲載の設定（JavaScript 用） ------------------------
\definecolor{cmtgreen}{rgb}{0.0,0.5,0.0}
\definecolor{kwblue}{rgb}{0.0,0.0,0.7}
\definecolor{strred}{rgb}{0.6,0.1,0.1}
\definecolor{bggray}{rgb}{0.97,0.97,0.97}

\lstdefinelanguage{JavaScriptCustom}{
  keywords={const,let,var,function,return,if,else,for,while,in,of,new,
            require,module,exports,true,false,null,undefined,async,await},
  ndkeywords={app,res,req,express,multer},
  sensitive=true,
  comment=[l]{//},
  morecomment=[s]{/*}{*/},
  string=[b]",
  morestring=[b]',
  morestring=[b]`,
}

\lstset{
  language=JavaScriptCustom,
  basicstyle=\ttfamily\scriptsize,
  keywordstyle=\color{kwblue}\bfseries,
  commentstyle=\color{cmtgreen}\itshape,
  stringstyle=\color{strred},
  backgroundcolor=\color{bggray},
  numbers=left,
  numberstyle=\tiny\color{gray},
  numbersep=8pt,
  frame=single,
  rulecolor=\color{gray},
  breaklines=true,
  breakatwhitespace=false,
  showstringspaces=false,
  columns=fullflexible,
  keepspaces=true,
  tabsize=2,
}

\begin{document}

\begin{center}
  {\LARGE\bfseries クラウドコンピューティング中間レポート}\\[4mm]
  {\large Node によるサーバー \& クライアントアプリケーションの制作}
\end{center}

\vspace{6mm}
\begin{flushright}
  \begin{tabular}{ll}
    学生番号 : & \studentID \\
    氏\hspace{1em}名 : & \studentName \\
    提出日\hspace{0.2em} : & 2026 年 6 月 9 日 \\
  \end{tabular}
\end{flushright}

\vspace{4mm}
\hrule
\vspace{6mm}

\section{目的}

Node.js と Express を用いて，\textbf{複数の自作 Web サイトを 1 つのサーバーで束ねる
「制作サイトまとめ」アプリケーション}を作ることを目的とした．
トップページから各サイトへ移動でき，それぞれのサイトで一覧・詳細・登録・編集・削除
といった一連の操作（CRUD）が行えることを目指した．具体的には次の機能を実装した．

\begin{itemize}
  \item \textbf{まとめトップページ}（\texttt{/}）：制作した各サイトへカード形式で移動できる
        入口ページ．
  \item \textbf{哲学者一覧サイト}（\texttt{/manga}）：漫画に登場する哲学者を題材にした
        一覧・詳細・新規登録・編集・削除．
  \item \textbf{教育学者一覧サイト}（\texttt{/kyouiku}）：主要な教育思想家の一覧・詳細・
        新規登録・編集・削除．
  \item \textbf{美少女一覧サイト}（\texttt{/bisyoujo}）：好きなキャラクターの一覧・詳細・
        新規登録・編集・削除．
  \item \textbf{画像アップロード機能}：新規登録・編集画面から各人物の画像（JPG）を
        アップロードでき，詳細ページに表示される．
  \item \textbf{おすすめ診断機能}（\texttt{/sindan}）：いくつかの質問に答えると，3 種類の
        テーマ（哲学者・教育学者・美少女）それぞれで「自分に合う一人」を判定し，
        相性ランキングを表示する診断サイト．
\end{itemize}

加えて，全ページのデザインを統一し（グラデーション背景・カード型 UI・ボタン形式の
ナビゲーション），使いやすく見栄えのする Web アプリケーションにすることを狙った．

% ===================== アクセスリンク =====================
\section{アクセスできるリンク}

下記 URL からブラウザでアクセスできる．

\begin{quote}
  \large \url{\siteURL}
\end{quote}

トップページから各サイト（哲学者・教育学者・美少女・おすすめ診断）へ移動できる．

% ===================== オリジナリティ =====================
\section{オリジナリティについて}

本課題は\textbf{すべて個人で制作したオリジナル作品}である．
題材の選定，扱うデータ（各人物の情報），ページ構成，デザイン，診断ロジック，
画像アップロード機能などをすべて自分で設計・実装した．
共同開発ではなく，クライアントサイド・サーバーサイドともに自分一人で担当した．

% ===================== ソースコード =====================
\section{サーバー側プログラム（ソースコード）}

サーバー側のプログラム \texttt{app5\_ver3.js} を以下に添付する．
処理の要所には日本語でコメントを記している．

\lstinputlisting[language=JavaScriptCustom]{app5_ver3.js}

\end{document}
